\section*{Abstract}

Troubleshooting computer networks is a complex, mostly manual activity. A network operator
works from the first symptom through device after device until a root cause is found. This
project investigates whether an AI-driven multi-agent system can take over this troubleshooting work
and at what benefit and cost. A system was built in which an orchestrator agent splits a
troubleshooting request and delegates the tasks to specialised agents, all operating against live network devices.

The system was evaluated in a containerlab environment of Nokia SR Linux devices using ten reproducible 
scenarios (a healthy baseline plus nine injected faults across three
difficulty tiers) and six language models from three price tiers. Each of the 300 runs
(10 scenarios $\times$ 6 models $\times$ 5 repetitions) was evaluated manually against the
scenario's known root cause.

The results show that the outcome depends mainly on the model used. With identical
agent logic, the success rate spans 16\,\% to 100\,\%. The high-tier models diagnosed 98 of 100
runs correctly across all difficulty tiers. They were typically within about two to three minutes and almost
independently of fault difficulty. Mid-tier models were reliable on easy faults (74--82\,\%
overall) but failed on harder faults, mostly
by hitting the time cap rather than by giving a wrong result. Low-tier models were
unusable and one of them repeatedly invented faults in a healthy network. With the mid- and high-tier
models, a correct diagnosis cost \$0.21--\$2.91 in API fees, a fraction of the
network operator time it replaces. Within the limits of a small single-fault lab, the experiments
show that autonomous multi-agent network troubleshooting works today when a sufficiently capable
model is used.

\vspace{2ex}
